% =========================================================
% Special Chapter: Mechanics Exam Practice
% POSN 1 Physics Preparation Notes
% Language: Thai
% Compiler: XeLaTeX
% =========================================================

\chapter{ชุดฝึกข้อสอบรวมกลศาสตร์}

\begin{center}
    {\Large \textbf{Mechanics Exam Practice}}\\[4pt]
    {\normalsize เอกสารเตรียมสอบคัดเลือก สอวน. ฟิสิกส์ ค่าย 1}\\[6pt]
    {\normalsize จัดทำโดย \textbf{Tames Thanakrit Damduan}}\\[2pt]
    {\footnotesize \textcopyright\ 2026 Tames Thanakrit Damduan. All rights reserved.}\\[8pt]
\end{center}

\section*{เป้าหมายของบทพิเศษ}

บทนี้ไม่ใช่บทเนื้อหาใหม่ แต่เป็นบทฝึกโจทย์รวมหลังเรียนกลศาสตร์ครบชุดแล้ว ได้แก่ การเคลื่อนที่แนวตรง เวกเตอร์ การเคลื่อนที่สองมิติ กฎของนิวตัน วงกลม งานและพลังงาน โมเมนตัม การหมุน โมเมนตัมเชิงมุม และของไหล

โจทย์ในบทนี้คัดเฉพาะข้อที่เกี่ยวข้องกับกลศาสตร์ และยังไม่ได้ถูกใส่ไว้ท้ายบทก่อนหน้า เพื่อไม่ให้ซ้ำกับชุดฝึกเดิม

\begin{tcolorbox}[title=แนวคิดในการทำโจทย์รวมกลศาสตร์]
โจทย์รวมกลศาสตร์มักไม่ได้บอกตรง ๆ ว่าต้องใช้บทใด นักเรียนต้องเริ่มจากวิเคราะห์สถานการณ์ก่อนว่าเป็นปัญหาแบบใด เช่น สมดุล แรงลัพธ์ พลังงาน โมเมนตัม การหมุน หรือของไหล จากนั้นค่อยเลือกเครื่องมือที่เหมาะสม ไม่ใช่เลือกสูตรจากคำสำคัญเพียงอย่างเดียว
\end{tcolorbox}

\section{แผนเลือกวิธีทำโจทย์}

\begin{tcolorbox}[title=แผนตัดสินใจเร็ว]
\begin{enumerate}
    \item ถ้าโจทย์ถามตำแหน่ง ความเร็ว เวลา หรือความสูง ให้เริ่มจากสมการการเคลื่อนที่
    \item ถ้าโจทย์มีแรงหลายแรง ให้วาด free-body diagram แล้วใช้
    \[
    \sum F=ma
    \]
    \item ถ้าโจทย์มีงาน ระยะทาง ความสูง สปริง หรือความเร็วที่ตำแหน่งต่าง ๆ ให้คิดด้วยพลังงาน
    \item ถ้าโจทย์มีการชน ระเบิด แรงเฉลี่ย หรือช่วงเวลาสั้น ๆ ให้คิดด้วยโมเมนตัมและอิมพัลส์
    \item ถ้าโจทย์มีคาน แท่ง ล้อ หรือวัตถุหมุน ให้คิดทอร์กและโมเมนต์ความเฉื่อย
    \item ถ้าโจทย์มีของเหลว ความดัน แรงพยุง หรือการไหล ให้แยกว่าเป็นของไหลอยู่นิ่งหรือกำลังไหล
\end{enumerate}
\end{tcolorbox}

% =========================================================
\section{ข้อสอบจริง สอวน. เพิ่มเติมที่ยังไม่ได้แทรกท้ายบท}

% =========================================================
\subsection{ปี 2560 ตอนที่ 2 ข้อ 1: มุมยิงเมื่อความสูงเป็นครึ่งหนึ่งของความสูงสูงสุด}

\vspace{1.5em}

\IfFileExists{img/posn60s2q1.png}{
\begin{center}
    \includegraphics[width=1\linewidth]{img/posn60s2q1.png}
\end{center}
}{
\begin{tcolorbox}[title=ตำแหน่งภาพที่แนะนำ]
ให้ตัดภาพข้อสอบปี 2560 ตอนที่ 2 ข้อ 1 แล้วบันทึกเป็น

\[
\texttt{img/posn60s2q1.png}
\]
\end{tcolorbox}
}

\begin{examplebox}{เฉลยละเอียด}
ให้ก้อนหินถูกขว้างด้วยความเร็วต้น

\[
u
\]

ทำมุม

\[
\theta
\]

กับแนวระดับ

องค์ประกอบความเร็วต้นคือ

\[
u_x=u\cos\theta
\]

\[
u_y=u\sin\theta
\]

ความสูงสูงสุดคือ

\[
H=\frac{u_y^2}{2g}
\]

โจทย์บอกว่าที่ความสูง

\[
\frac{H}{2}
\]

ทิศของความเร็วทำมุม $60^\circ$ กับแนวระดับ

ที่ความสูง $H/2$ ใช้สมการแนวดิ่ง

\[
v_y^2=u_y^2-2g\left(\frac{H}{2}\right)
\]

แทน

\[
H=\frac{u_y^2}{2g}
\]

จะได้

\[
v_y^2=u_y^2-gH
\]

\[
v_y^2=u_y^2-\frac{u_y^2}{2}
\]

\[
v_y^2=\frac{u_y^2}{2}
\]

ดังนั้น

\[
v_y=\frac{u_y}{\sqrt{2}}
\]

ความเร็วแนวราบคงตัว

\[
v_x=u_x
\]

จากเงื่อนไขมุมของความเร็ว

\[
\tan60^\circ=\frac{v_y}{v_x}
\]

แทนค่า

\[
\sqrt{3}=
\frac{\frac{u_y}{\sqrt{2}}}{u_x}
\]

\[
\sqrt{3}=
\frac{u\sin\theta}{\sqrt{2}u\cos\theta}
\]

\[
\sqrt{3}=\frac{\tan\theta}{\sqrt{2}}
\]

ดังนั้น

\[
\tan\theta=\sqrt{6}
\]

สรุป

\[
\boxed{\theta=\tan^{-1}\sqrt{6}}
\]
\end{examplebox}

% =========================================================
\subsection{ปี 2560 ตอนที่ 2 ข้อ 2: ดันมวลใหญ่ให้มวลเล็กไม่ไถลลง}

\vspace{1.5em}

\IfFileExists{img/posn60s2q2.png}{
\begin{center}
    \includegraphics[width=1\linewidth]{img/posn60s2q2.png}
\end{center}
}{
\begin{tcolorbox}[title=ตำแหน่งภาพที่แนะนำ]
ให้ตัดภาพข้อสอบปี 2560 ตอนที่ 2 ข้อ 2 แล้วบันทึกเป็น

\[
\texttt{img/posn60s2q2.png}
\]
\end{tcolorbox}
}

\begin{examplebox}{เฉลยละเอียด}
มวลเล็ก $m$ ติดอยู่ด้านข้างของมวลใหญ่ $M$ และพื้นระดับลื่น ต้องออกแรง $F$ ให้ระบบเร่งในแนวราบมากพอ เพื่อให้แรงเสียดทานพยุงมวลเล็กไม่ให้ไถลลง

ถ้ามวลทั้งสองเคลื่อนที่ไปด้วยกัน ความเร่งแนวราบของระบบคือ

\[
a=\frac{F}{M+m}
\]

พิจารณามวลเล็ก $m$

แรงปกติจากผนังของมวลใหญ่เป็นแรงแนวราบ และเป็นแรงที่ทำให้มวลเล็กเร่ง

\[
N=ma
\]

แรงเสียดทานสถิตสูงสุดคือ

\[
f_{s,\max}=\mu_sN
\]

เพื่อไม่ให้มวลเล็กไถลลง ต้องมี

\[
f_{s,\max}\geq mg
\]

ดังนั้น

\[
\mu_sN\geq mg
\]

แทน

\[
N=ma
\]

จะได้

\[
\mu_sma\geq mg
\]

ตัด $m$ ออก

\[
a\geq \frac{g}{\mu_s}
\]

เพราะ

\[
a=\frac{F}{M+m}
\]

ดังนั้น

\[
\frac{F}{M+m}\geq \frac{g}{\mu_s}
\]

จึงได้

\[
F_{\min}=\frac{(M+m)g}{\mu_s}
\]

โจทย์ให้

\[
\mu_s=0.50
\]

ดังนั้น

\[
\boxed{F_{\min}=2(M+m)g}
\]
\end{examplebox}

% =========================================================
\subsection{ปี 2560 ตอนที่ 2 ข้อ 3: ท่อนวัตถุไม่สม่ำเสมอและสมดุลทอร์ก}

\vspace{1.5em}

\IfFileExists{img/posn60s2q3.png}{
\begin{center}
    \includegraphics[width=1\linewidth]{img/posn60s2q3.png}
\end{center}
}{
\begin{tcolorbox}[title=ตำแหน่งภาพที่แนะนำ]
ให้ตัดภาพข้อสอบปี 2560 ตอนที่ 2 ข้อ 3 แล้วบันทึกเป็น

\[
\texttt{img/posn60s2q3.png}
\]
\end{tcolorbox}
}

\begin{examplebox}{เฉลยละเอียด}
ให้ท่อนวัตถุยาว

\[
3L
\]

มีน้ำหนักรวม

\[
W
\]

และจุดศูนย์กลางมวลอยู่ห่างจากปลาย $A$ เป็นระยะ

\[
x
\]

\textbf{กรณีแรก}

มีเชือกดึงที่ปลาย $A$ และปลาย $B$ โดยแรงตึงที่ปลาย $A$ เท่ากับ

\[
T
\]

ดังนั้นแรงตึงที่ปลาย $B$ คือ

\[
W-T
\]

ตั้งสมดุลทอร์กรอบปลาย $A$

\[
(W-T)(3L)=Wx
\]

\[
Wx=3L(W-T)
\]

\textbf{กรณีที่สอง}

เอาเชือกที่ปลาย $A$ ออก แล้วใช้ลิ่มค้ำที่ระยะ

\[
L
\]

จากปลาย $A$ และแรงตึงเชือกที่ปลาย $B$ ยังเท่ากับ

\[
T
\]

ให้แรงจากลิ่มเป็น

\[
R
\]

สมดุลแรงแนวดิ่งให้

\[
R+T=W
\]

ดังนั้น

\[
R=W-T
\]

ตั้งสมดุลทอร์กรอบปลาย $A$

\[
R(L)+T(3L)=Wx
\]

แทน $R=W-T$

\[
(W-T)L+3TL=Wx
\]

แต่จากกรณีแรก

\[
Wx=3L(W-T)
\]

ดังนั้น

\[
(W-T)L+3TL=3L(W-T)
\]

ตัด $L$ ออก

\[
W-T+3T=3W-3T
\]

\[
W+2T=3W-3T
\]

\[
5T=2W
\]

ดังนั้น

\[
W=\frac{5}{2}T
\]

สรุป

\[
\boxed{\frac{W}{T}=\frac{5}{2}}
\]
\end{examplebox}

% =========================================================
\subsection{ปี 2560 ข้อ 10: คานสม่ำเสมอ เชือกเอียง และน้ำหนักแขวน}

\vspace{1.5em}

\IfFileExists{img/posn60q10.png}{
\begin{center}
    \includegraphics[width=1\linewidth]{img/posn60q10.png}
\end{center}
}{
\begin{tcolorbox}[title=ตำแหน่งภาพที่แนะนำ]
ให้ตัดภาพข้อสอบปี 2560 ข้อ 10 แล้วบันทึกเป็น

\[
\texttt{img/posn60q10.png}
\]
\end{tcolorbox}
}

\begin{examplebox}{เฉลยละเอียด}
คานยาว

\[
4.0\,\si{m}
\]

มีน้ำหนัก

\[
200\,\si{N}
\]

กระทำที่กึ่งกลางคาน ห่างจากบานพับ

\[
2.0\,\si{m}
\]

ที่ปลายคานมีน้ำหนักแขวน

\[
300\,\si{N}
\]

ห่างจากบานพับ

\[
4.0\,\si{m}
\]

เชือกจากปลายคานไปยังกำแพงเป็นสามเหลี่ยม $3-4-5$ ดังนั้นองค์ประกอบแนวดิ่งของแรงตึงคือ

\[
T_y=T\sin\phi=T\left(\frac{3}{5}\right)
\]

ตั้งสมดุลทอร์กรอบบานพับ

\[
T_y(4.0)=200(2.0)+300(4.0)
\]

แทน $T_y=\frac{3}{5}T$

\[
\left(\frac{3}{5}T\right)(4.0)=400+1200
\]

\[
\frac{12}{5}T=1600
\]

\[
T=\frac{1600(5)}{12}
\]

\[
T\approx 667\,\si{N}
\]

ดังนั้น

\[
\boxed{T\approx 667\,\si{N}}
\]
\end{examplebox}

% =========================================================
\subsection{ปี 2561 ข้อ 8: โพรเจคไทล์ชนลูกโป่งที่เร่งขึ้น}

\vspace{1.5em}

\IfFileExists{img/posn61q8.png}{
\begin{center}
    \includegraphics[width=1\linewidth]{img/posn61q8.png}
\end{center}
}{
\begin{tcolorbox}[title=ตำแหน่งภาพที่แนะนำ]
ให้ตัดภาพข้อสอบปี 2561 ข้อ 8 แล้วบันทึกเป็น

\[
\texttt{img/posn61q8.png}
\]
\end{tcolorbox}
}

\begin{examplebox}{เฉลยละเอียด}
วัตถุถูกขว้างด้วยความเร็วต้น

\[
u=30\,\si{m/s}
\]

ทำมุม

\[
30^\circ
\]

กับแนวระดับ

ระยะห่างแนวราบถึงลูกโป่งคือ

\[
30\sqrt{3}\,\si{m}
\]

หาเวลาไปถึงแนวของลูกโป่งจากแกน $x$

\[
x=u\cos30^\circ t
\]

ดังนั้น

\[
30\sqrt{3}=30\left(\frac{\sqrt{3}}{2}\right)t
\]

\[
t=2\,\si{s}
\]

ตำแหน่งแนวดิ่งของโพรเจคไทล์ที่เวลานี้คือ

\[
y=u\sin30^\circ t-\frac{1}{2}gt^2
\]

ใช้ $g=9.8\,\si{m/s^2}$

\[
y=30\left(\frac{1}{2}\right)(2)-\frac{1}{2}(9.8)(2^2)
\]

\[
y=30-19.6
\]

\[
y=10.4\,\si{m}
\]

ลูกโป่งเริ่มจากหยุดนิ่งและเร่งขึ้นด้วยความเร่ง $a$ ดังนั้น

\[
y_{\text{balloon}}=\frac{1}{2}at^2
\]

เพื่อให้ชนกัน

\[
\frac{1}{2}a(2^2)=10.4
\]

\[
2a=10.4
\]

\[
a=5.2\,\si{m/s^2}
\]

ดังนั้น

\[
\boxed{a_{\min}=5.2\,\si{m/s^2}}
\]
\end{examplebox}

% =========================================================
\subsection{ปี 2561 ข้อ 9: ลูกตุ้มกรวยสองลูกและอัตราส่วนแรงตึงเชือก}

\vspace{1.5em}

\IfFileExists{img/posn61q9.png}{
\begin{center}
    \includegraphics[width=1\linewidth]{img/posn61q9.png}
\end{center}
}{
\begin{tcolorbox}[title=ตำแหน่งภาพที่แนะนำ]
ให้ตัดภาพข้อสอบปี 2561 ข้อ 9 แล้วบันทึกเป็น

\[
\texttt{img/posn61q9.png}
\]
\end{tcolorbox}
}

\begin{examplebox}{เฉลยละเอียด}
ให้มวลของลูก A เป็นสองเท่าของลูก B

\[
m_A=2m_B
\]

เชือกของลูก B ยาวเป็น $1.5$ เท่าของเชือก A

\[
L_B=\frac{3}{2}L_A
\]

คาบของ A เป็นสองเท่าของ B

\[
T_A=2T_B
\]

สำหรับลูกตุ้มกรวย

\[
T_{\text{period}}=2\pi\sqrt{\frac{L\cos\theta}{g}}
\]

ดังนั้น

\[
T_{\text{period}}^2\propto L\cos\theta
\]

จึงได้

\[
\frac{T_A^2}{T_B^2}
=
\frac{L_A\cos\theta_A}{L_B\cos\theta_B}
\]

แทนค่า

\[
4=
\frac{L_A\cos\theta_A}{\frac{3}{2}L_A\cos\theta_B}
\]

\[
4=\frac{2}{3}\frac{\cos\theta_A}{\cos\theta_B}
\]

ดังนั้น

\[
\frac{\cos\theta_A}{\cos\theta_B}=6
\]

แรงตึงเชือกของลูกตุ้มกรวยได้จากสมดุลแนวดิ่ง

\[
T\cos\theta=mg
\]

ดังนั้น

\[
T=\frac{mg}{\cos\theta}
\]

หาอัตราส่วน

\[
\frac{T_A}{T_B}
=
\frac{\frac{m_Ag}{\cos\theta_A}}{\frac{m_Bg}{\cos\theta_B}}
\]

\[
\frac{T_A}{T_B}
=
\frac{m_A}{m_B}
\frac{\cos\theta_B}{\cos\theta_A}
\]

แทน

\[
\frac{m_A}{m_B}=2
\]

และ

\[
\frac{\cos\theta_B}{\cos\theta_A}=\frac{1}{6}
\]

จะได้

\[
\frac{T_A}{T_B}=2\left(\frac{1}{6}\right)
\]

\[
\frac{T_A}{T_B}=\frac{1}{3}
\]

สรุป

\[
\boxed{\frac{T_A}{T_B}=\frac{1}{3}}
\]
\end{examplebox}

% =========================================================
\subsection{ปี 2561 ตอนที่ 2 ข้อ 3: ทรงกระบอกบนพื้นเอียงและเชือกแนวระดับ}

\vspace{1.5em}

\IfFileExists{img/posn61s2q3.png}{
\begin{center}
    \includegraphics[width=1\linewidth]{img/posn61s2q3.png}
\end{center}
}{
\begin{tcolorbox}[title=ตำแหน่งภาพที่แนะนำ]
ให้ตัดภาพข้อสอบปี 2561 ตอนที่ 2 ข้อ 3 แล้วบันทึกเป็น

\[
\texttt{img/posn61s2q3.png}
\]
\end{tcolorbox}
}

\begin{examplebox}{เฉลยละเอียด}
ทรงกระบอกอยู่ในสมดุลบนพื้นเอียงมุม

\[
60^\circ
\]

ให้เชือกดึงในแนวระดับ และให้แรงเสียดทานที่พื้นเอียงมีค่า

\[
f
\]

น้ำหนักของทรงกระบอกคือ

\[
W
\]

พิจารณาทอร์กรอบจุดศูนย์กลางของทรงกระบอก

แรงปกติและน้ำหนักผ่านแนวศูนย์กลาง จึงไม่ให้ทอร์กรอบศูนย์กลาง

แรงตึงเชือกและแรงเสียดทานให้ทอร์กตรงข้ามกัน และมีระยะแขนเท่ากับรัศมี $R$ ทั้งคู่

ดังนั้นสมดุลทอร์กให้

\[
TR=fR
\]

จึงได้

\[
f=T
\]

พิจารณาสมดุลแรงตามแนวพื้นเอียง

องค์ประกอบของแรงตึงตามแนวพื้นเอียงคือ

\[
T\cos60^\circ
\]

แรงเสียดทานมีทิศขึ้นตามพื้นเอียง และมีขนาด $T$

องค์ประกอบของน้ำหนักลงตามพื้นเอียงคือ

\[
W\sin60^\circ
\]

สมดุลตามแนวพื้นเอียง

\[
T\cos60^\circ+T=W\sin60^\circ
\]

แทน

\[
\cos60^\circ=\frac{1}{2}
\]

และ

\[
\sin60^\circ=\frac{\sqrt{3}}{2}
\]

จะได้

\[
\frac{1}{2}T+T=\frac{\sqrt{3}}{2}W
\]

\[
\frac{3}{2}T=\frac{\sqrt{3}}{2}W
\]

ดังนั้น

\[
T=\frac{W}{\sqrt{3}}
\]

สรุป

\[
\boxed{T=\frac{W}{\sqrt{3}}}
\]
\end{examplebox}

% =========================================================
\subsection{ปี 2561 ตอนที่ 2 ข้อ 6: ดึงวัตถุขึ้นในแนวดิ่งด้วยแรงคงตัว}

\vspace{1.5em}

\IfFileExists{img/posn61s2q6.png}{
\begin{center}
    \includegraphics[width=1\linewidth]{img/posn61s2q6.png}
\end{center}
}{
\begin{tcolorbox}[title=ตำแหน่งภาพที่แนะนำ]
ให้ตัดภาพข้อสอบปี 2561 ตอนที่ 2 ข้อ 6 แล้วบันทึกเป็น

\[
\texttt{img/posn61s2q6.png}
\]
\end{tcolorbox}
}

\begin{examplebox}{เฉลยละเอียด}
วัตถุมวล

\[
m=1.0\,\si{kg}
\]

ถูกดึงขึ้นในแนวดิ่งจากหยุดนิ่งด้วยแรงคงตัว

\[
F=20\,\si{N}
\]

เป็นระยะทาง

\[
s=5.0\,\si{m}
\]

ต้องการหาการเพิ่มขึ้นของพลังงานจลน์

ใช้ทฤษฎีงานกับพลังงานจลน์

\[
\Delta K=W_{\text{net}}
\]

งานของแรงดึงคือ

\[
W_F=Fs
\]

งานของน้ำหนักเป็นลบ

\[
W_g=-mgs
\]

ดังนั้น

\[
\Delta K=Fs-mgs
\]

แทนค่า

\[
\Delta K=(20)(5.0)-(1.0)(9.8)(5.0)
\]

\[
\Delta K=100-49
\]

\[
\Delta K=51\,\si{J}
\]

ดังนั้น

\[
\boxed{\Delta K=51\,\si{J}}
\]
\end{examplebox}

% =========================================================
\subsection{ปี 2562 ข้อ 12: เชือกหนักห้อยท้องช้าง}

\vspace{1.5em}

\IfFileExists{img/posn62q12.png}{
\begin{center}
    \includegraphics[width=1\linewidth]{img/posn62q12.png}
\end{center}
}{
\begin{tcolorbox}[title=ตำแหน่งภาพที่แนะนำ]
ให้ตัดภาพข้อสอบปี 2562 ข้อ 12 แล้วบันทึกเป็น

\[
\texttt{img/posn62q12.png}
\]
\end{tcolorbox}
}

\begin{examplebox}{เฉลยละเอียด}
เชือกหนักรวม

\[
W
\]

ห้อยระหว่างจุดแขวนสองจุดที่อยู่ระดับเดียวกัน และเชือกที่จุดแขวนทำมุม

\[
\theta
\]

กับแนวดิ่ง

ใช้สมมาตร พิจารณาครึ่งหนึ่งของเชือก

น้ำหนักของครึ่งเชือกคือ

\[
\frac{W}{2}
\]

ให้แรงตึงที่จุดแขวนเป็น

\[
T_s
\]

แรงตึงที่จุดต่ำสุดเป็นแนวนอน ให้เป็น

\[
T_0
\]

สมดุลแนวดิ่งของครึ่งเชือก

\[
T_s\cos\theta=\frac{W}{2}
\]

สมดุลแนวนอน

\[
T_0=T_s\sin\theta
\]

จากสมการแรก

\[
T_s=\frac{W}{2\cos\theta}
\]

ดังนั้น

\[
T_0=\frac{W}{2\cos\theta}\sin\theta
\]

\[
T_0=\frac{W}{2}\tan\theta
\]

สรุป

\[
\boxed{T_0=\frac{W}{2}\tan\theta}
\]
\end{examplebox}

% =========================================================
\subsection{ปี 2562 ข้อ 15: เหรียญกลมชนเยื้องศูนย์}

\vspace{1.5em}

\IfFileExists{img/posn62q15.png}{
\begin{center}
    \includegraphics[width=1\linewidth]{img/posn62q15.png}
\end{center}
}{
\begin{tcolorbox}[title=ตำแหน่งภาพที่แนะนำ]
ให้ตัดภาพข้อสอบปี 2562 ข้อ 15 แล้วบันทึกเป็น

\[
\texttt{img/posn62q15.png}
\]
\end{tcolorbox}
}

\begin{examplebox}{เฉลยละเอียด}
เหรียญสองเหรียญมีรัศมีเท่ากันคือ

\[
R
\]

และขอบเกลี้ยง เมื่อชนกัน แรงระหว่างเหรียญอยู่ตามแนวเส้นเชื่อมศูนย์กลางของเหรียญทั้งสอง

ดังนั้นทิศของความเร็วของเหรียญ $B$ หลังชนจะอยู่ตามแนวเส้นเชื่อมศูนย์กลางขณะชน

จากรูป ศูนย์กลางของเหรียญ $B$ สูงกว่าศูนย์กลางของเหรียญ $A$ เป็นระยะ

\[
a
\]

ระยะระหว่างศูนย์กลางขณะสัมผัสกันคือ

\[
2R
\]

ดังนั้น ถ้าเส้นเชื่อมศูนย์กลางทำมุม $\theta$ กับแนวระดับ จะมี

\[
\sin\theta=\frac{a}{2R}
\]

จึงได้

\[
\boxed{\theta=\sin^{-1}\left(\frac{a}{2R}\right)}
\]
\end{examplebox}

% =========================================================
\subsection{ปี 2562 ข้อ 24: ลูกตุ้มสองลูกชนกันซ้ำ}

\vspace{1.5em}

\IfFileExists{img/posn62q24.png}{
\begin{center}
    \includegraphics[width=1\linewidth]{img/posn62q24.png}
\end{center}
}{
\begin{tcolorbox}[title=ตำแหน่งภาพที่แนะนำ]
ให้ตัดภาพข้อสอบปี 2562 ข้อ 24 แล้วบันทึกเป็น

\[
\texttt{img/posn62q24.png}
\]
\end{tcolorbox}
}

\begin{examplebox}{เฉลยแนวคิด}
ลูกตุ้มสองลูกมีความยาวเชือกเท่ากัน

\[
\ell
\]

ดังนั้นคาบของการแกว่งเล็ก ๆ เท่ากัน

\[
T=2\pi\sqrt{\frac{\ell}{g}}
\]

หลังชนครั้งแรก ลูกตุ้มทั้งสองแยกออกจากตำแหน่งล่างสุด แล้วจะกลับมาพบกันอีกครั้งที่ตำแหน่งล่างสุดหลังเวลาครึ่งคาบ

ดังนั้นเวลาระหว่างการชนครั้งแรกกับครั้งที่สองคือ

\[
\Delta t=\frac{T}{2}
\]

แทนค่า $T$

\[
\Delta t=\pi\sqrt{\frac{\ell}{g}}
\]

สรุป

\[
\boxed{\Delta t=\pi\sqrt{\frac{\ell}{g}}}
\]
\end{examplebox}

% =========================================================
\subsection{ปี 2562 ข้อ 26: เหรียญสามเหรียญชนกันแบบยืดหยุ่นสมมาตร}

\vspace{1.5em}

\IfFileExists{img/posn62q26.png}{
\begin{center}
    \includegraphics[width=1\linewidth]{img/posn62q26.png}
\end{center}
}{
\begin{tcolorbox}[title=ตำแหน่งภาพที่แนะนำ]
ให้ตัดภาพข้อสอบปี 2562 ข้อ 26 แล้วบันทึกเป็น

\[
\texttt{img/posn62q26.png}
\]
\end{tcolorbox}
}

\begin{examplebox}{เฉลยละเอียด}
เหรียญทั้งสามมีมวลเท่ากัน และการชนเป็นแบบยืดหยุ่นสมมาตร

ก่อนชน เหรียญ $A$ มีอัตราเร็ว

\[
u
\]

หลังชน ให้เหรียญ $A$ สะท้อนกลับด้วยอัตราเร็ว

\[
v
\]

และเหรียญ $B,C$ เคลื่อนที่ออกแบบสมมาตรด้วยอัตราเร็วเท่ากัน

\[
w
\]

จากเรขาคณิตของเหรียญเท่ากันสามเหรียญ แนวการเคลื่อนที่ของ $B$ และ $C$ ทำมุม

\[
30^\circ
\]

กับแนวเดิมของ $A$

สมการโมเมนตัมตามแกน $x$

\[
mu=-mv+2mw\cos30^\circ
\]

ตัด $m$

\[
u=-v+2w\cos30^\circ
\]

เพราะ

\[
\cos30^\circ=\frac{\sqrt{3}}{2}
\]

จึงได้

\[
u=-v+\sqrt{3}w
\]

ดังนั้น

\[
v=\sqrt{3}w-u
\]

สมการพลังงานจลน์ เพราะชนยืดหยุ่น

\[
\frac{1}{2}mu^2=\frac{1}{2}mv^2+2\left(\frac{1}{2}mw^2\right)
\]

ตัด $\frac{1}{2}m$ ออก

\[
u^2=v^2+2w^2
\]

แทน

\[
v=\sqrt{3}w-u
\]

\[
u^2=(\sqrt{3}w-u)^2+2w^2
\]

\[
u^2=3w^2-2\sqrt{3}uw+u^2+2w^2
\]

ตัด $u^2$ ทั้งสองข้าง

\[
0=5w^2-2\sqrt{3}uw
\]

ไม่เอาคำตอบ $w=0$ จึงได้

\[
w=\frac{2\sqrt{3}}{5}u
\]

หา $v$

\[
v=\sqrt{3}\left(\frac{2\sqrt{3}}{5}u\right)-u
\]

\[
v=\frac{6}{5}u-u
\]

\[
v=\frac{1}{5}u
\]

สรุป

\[
\boxed{v=\frac{u}{5}}
\]

\[
\boxed{w=\frac{2\sqrt{3}}{5}u}
\]
\end{examplebox}

% =========================================================
\subsection{ปี 2566 ข้อ 2: ท่อนโตสม่ำเสมอหลังดึงตัวค้ำออก}

\vspace{1.5em}

\IfFileExists{img/posn66q2.png}{
\begin{center}
    \includegraphics[width=1\linewidth]{img/posn66q2.png}
\end{center}
}{
\begin{tcolorbox}[title=ตำแหน่งภาพที่แนะนำ]
ให้ตัดภาพข้อสอบปี 2566 ข้อ 2 แล้วบันทึกเป็น

\[
\texttt{img/posn66q2.png}
\]
\end{tcolorbox}
}

\begin{examplebox}{เฉลยละเอียด}
ท่อนโตสม่ำเสมอมวล

\[
M
\]

ยาว

\[
L
\]

ถูกค้ำที่จุด $A$ และ $B$ โดยจุด $A$ อยู่ห่างจากปลายซ้าย

\[
\frac{L}{4}
\]

ดังนั้นจุดศูนย์กลางมวลอยู่ห่างจากจุด $A$

\[
\frac{L}{4}
\]

เมื่อดึงตัวค้ำที่ $B$ ออกทันที ท่อนโตเริ่มหมุนรอบจุด $A$

โมเมนต์ความเฉื่อยรอบจุด $A$ คือ

\[
I_A=I_{\text{cm}}+M\left(\frac{L}{4}\right)^2
\]

สำหรับแท่งสม่ำเสมอ

\[
I_{\text{cm}}=\frac{1}{12}ML^2
\]

ดังนั้น

\[
I_A=\frac{1}{12}ML^2+\frac{1}{16}ML^2
\]

\[
I_A=\frac{7}{48}ML^2
\]

ทอร์กของน้ำหนักรอบจุด $A$ คือ

\[
\tau=Mg\left(\frac{L}{4}\right)
\]

ใช้สมการการหมุน

\[
\tau=I_A\alpha
\]

\[
Mg\left(\frac{L}{4}\right)=\frac{7}{48}ML^2\alpha
\]

\[
\alpha=\frac{12g}{7L}
\]

ความเร่งของจุดศูนย์กลางมวลทันทีหลังปล่อยคือ

\[
a_{\text{cm}}=\alpha\left(\frac{L}{4}\right)
\]

\[
a_{\text{cm}}=\frac{12g}{7L}\left(\frac{L}{4}\right)
\]

\[
a_{\text{cm}}=\frac{3g}{7}
\]

ทิศลง

ให้แรงที่ตัวค้ำ $A$ กระทำต่อแท่งเป็น $N$

เลือกลงเป็นบวก

\[
Mg-N=Ma_{\text{cm}}
\]

\[
Mg-N=M\left(\frac{3g}{7}\right)
\]

\[
N=\frac{4}{7}Mg
\]

ดังนั้นแรงที่แท่งกระทำต่อตัวค้ำ $A$ มีขนาดเท่ากัน

\[
\boxed{N=\frac{4}{7}Mg}
\]
\end{examplebox}

% =========================================================
\subsection{ปี 2566 ข้อ 18: ระบบดาวสองดวงโคจรรอบจุดศูนย์กลางมวล}

\vspace{1.5em}

\IfFileExists{img/posn66q18.png}{
\begin{center}
    \includegraphics[width=1\linewidth]{img/posn66q18.png}
\end{center}
}{
\begin{tcolorbox}[title=ตำแหน่งภาพที่แนะนำ]
ให้ตัดภาพข้อสอบปี 2566 ข้อ 18 แล้วบันทึกเป็น

\[
\texttt{img/posn66q18.png}
\]
\end{tcolorbox}
}

\begin{examplebox}{เฉลยละเอียด}
ระบบดาวสองดวงมวล

\[
M
\]

และ

\[
m
\]

อยู่ห่างกันคงที่

\[
R
\]

และโคจรรอบจุดศูนย์กลางมวลร่วม

ให้ระยะจากดาวมวล $M$ ถึงจุดศูนย์กลางมวลเป็น

\[
r_M
\]

และระยะจากดาวมวล $m$ ถึงจุดศูนย์กลางมวลเป็น

\[
r_m
\]

เพราะระยะห่างระหว่างดาวคือ

\[
R
\]

จึงมี

\[
r_M+r_m=R
\]

จากนิยามจุดศูนย์กลางมวล

\[
Mr_M=mr_m
\]

ต้องการหา $r_m$

จากสมการแรก

\[
r_M=R-r_m
\]

แทนลงในสมการศูนย์กลางมวล

\[
M(R-r_m)=mr_m
\]

\[
MR-Mr_m=mr_m
\]

\[
MR=(M+m)r_m
\]

ดังนั้น

\[
r_m=\frac{M}{M+m}R
\]

สรุป

\[
\boxed{r_m=\frac{M}{M+m}R}
\]
\end{examplebox}

% =========================================================
\section{โจทย์จากเอกสารฝึกระดับ สอวน. เพิ่มเติม}

\subsection{เอกสารเพิ่มเติม ข้อ 11: ถังน้ำถูกกระชากขึ้นและความดันที่ก้นถัง}

\vspace{1.5em}

\IfFileExists{img/exam_extra_q11.png}{
\begin{center}
    \includegraphics[width=1\linewidth]{img/exam_extra_q11.png}
\end{center}
}{
\begin{tcolorbox}[title=ตำแหน่งภาพที่แนะนำ]
ให้ตัดภาพโจทย์จากเอกสารเพิ่มเติมข้อ 11 แล้วบันทึกเป็น

\[
\texttt{img/exam\_extra\_q11.png}
\]
\end{tcolorbox}
}

\begin{examplebox}{เฉลยละเอียด}
ถังน้ำถูกเร่งขึ้นในแนวดิ่งด้วยความเร่ง

\[
a
\]

พิจารณาในกรอบของถัง จะมีความเร่งโน้มถ่วงเสมือนเพิ่มขึ้นเป็น

\[
g_{\text{eff}}=g+a
\]

ดังนั้นความดันเกจที่ก้นถังลึกจากผิวน้ำระยะ

\[
h
\]

คือ

\[
P_{\text{gauge}}=\rho(g+a)h
\]

ถ้าผิวน้ำเปิดสู่อากาศ ความดันสัมบูรณ์ที่ก้นถังคือ

\[
P_{\text{bottom}}=P_a+\rho(g+a)h
\]

สรุป

\[
\boxed{P_{\text{bottom}}=P_a+\rho(g+a)h}
\]
\end{examplebox}

% =========================================================
\subsection{เอกสารเพิ่มเติม ข้อ 12: น้ำพุ่งจากรูด้านข้างของถัง}

\vspace{1.5em}

\IfFileExists{img/exam_extra_q12.png}{
\begin{center}
    \includegraphics[width=1\linewidth]{img/exam_extra_q12.png}
\end{center}
}{
\begin{tcolorbox}[title=ตำแหน่งภาพที่แนะนำ]
ให้ตัดภาพโจทย์จากเอกสารเพิ่มเติมข้อ 12 แล้วบันทึกเป็น

\[
\texttt{img/exam\_extra\_q12.png}
\]
\end{tcolorbox}
}

\begin{examplebox}{เฉลยละเอียด}
น้ำในถังสูงจากพื้นถึงผิวน้ำเป็น

\[
h
\]

รูอยู่สูงจากพื้น

\[
y
\]

ดังนั้นความลึกของรูใต้ผิวน้ำคือ

\[
h-y
\]

จากกฎของตอร์ริเชลลี อัตราเร็วของน้ำที่พุ่งออกจากรูคือ

\[
v=\sqrt{2g(h-y)}
\]

น้ำพุ่งออกแนวระดับและตกถึงพื้นที่ระดับต่ำกว่าเป็นระยะ

\[
y
\]

เวลาตกถึงพื้นหาได้จากแนวดิ่ง

\[
y=\frac{1}{2}gt^2
\]

ดังนั้น

\[
t=\sqrt{\frac{2y}{g}}
\]

ระยะทางแนวราบคือ

\[
x=vt
\]

แทนค่า

\[
x=\sqrt{2g(h-y)}\sqrt{\frac{2y}{g}}
\]

\[
x=2\sqrt{y(h-y)}
\]

ยกกำลังสอง

\[
x^2=4y(h-y)
\]

จัดรูป

\[
x^2=4hy-4y^2
\]

และเขียนเป็นรูปกำลังสองสมบูรณ์

\[
x^2=h^2-4\left(y-\frac{h}{2}\right)^2
\]

ดังนั้นพจน์ในช่องว่างคือ

\[
\boxed{y-\frac{h}{2}}
\]
\end{examplebox}


% =========================================================
% Additional POSN 2563 Mechanics Problems
% Append this block at the end of the Mechanics Exam Practice chapter
% =========================================================

% =========================================================
\subsection{ปี 2563 ข้อ 2: ลูกโป่งจมอยู่ในของเหลวและแรงตึงเชือก}

\vspace{1.5em}

\IfFileExists{img/posn63q2.png}{
\begin{center}
    \includegraphics[width=1\linewidth]{img/posn63q2.png}
\end{center}
}{
\begin{tcolorbox}[title=ตำแหน่งภาพที่แนะนำ]
ให้ตัดภาพข้อสอบปี 2563 ข้อ 2 แล้วบันทึกเป็น

\[
\texttt{img/posn63q2.png}
\]
\end{tcolorbox}
}

\begin{examplebox}{เฉลยละเอียด}
ลูกโป่งมวล $m$ และปริมาตร $V$ จมอยู่ในของเหลวความหนาแน่น $\rho$ โดยมีเชือกรั้งลงด้านล่าง

แรงที่กระทำต่อลูกโป่งมีสามแรง

\[
\text{แรงพยุง } B=\rho Vg \text{ ชี้ขึ้น}
\]

\[
\text{น้ำหนัก } mg \text{ ชี้ลง}
\]

\[
\text{แรงตึงเชือก } T \text{ ชี้ลง}
\]

เมื่อลูกโป่งอยู่นิ่ง แรงลัพธ์เป็นศูนย์

\[
B-mg-T=0
\]

แทน

\[
B=\rho Vg
\]

จะได้

\[
\rho Vg-mg-T=0
\]

ดังนั้น

\[
\boxed{T=(\rho V-m)g}
\]
\end{examplebox}

% =========================================================
\subsection{ปี 2563 ข้อ 3: การชนยืดหยุ่นของมวลเท่ากันและกำแพงแข็ง}

\vspace{1.5em}

\IfFileExists{img/posn63q3.png}{
\begin{center}
    \includegraphics[width=1\linewidth]{img/posn63q3.png}
\end{center}
}{
\begin{tcolorbox}[title=ตำแหน่งภาพที่แนะนำ]
ให้ตัดภาพข้อสอบปี 2563 ข้อ 3 แล้วบันทึกเป็น

\[
\texttt{img/posn63q3.png}
\]
\end{tcolorbox}
}

\begin{examplebox}{เฉลยละเอียด}
มวล $A$ และ $B$ เท่ากัน พื้นลื่น และการชนเป็นแบบยืดหยุ่น

ก่อนชนครั้งแรก

\[
A \text{ เคลื่อนที่ด้วยความเร็ว } v
\]

\[
B \text{ อยู่นิ่ง}
\]

สำหรับการชนยืดหยุ่นหนึ่งมิติของมวลเท่ากัน วัตถุทั้งสองจะแลกเปลี่ยนความเร็วกัน หลังชนครั้งแรกจึงได้

\[
v_A=0
\]

\[
v_B=v
\]

มวล $B$ เคลื่อนที่ไปชนกำแพงซึ่งอยู่ห่างออกไปเป็นระยะ $D$ ใช้เวลา

\[
t_1=\frac{D}{v}
\]

เมื่อชนกำแพงอย่างยืดหยุ่น $B$ สะท้อนกลับด้วยอัตราเร็วเดิม $v$ และเคลื่อนที่กลับมาหา $A$ ซึ่งหยุดอยู่ที่ตำแหน่งเดิม ใช้เวลาอีก

\[
t_2=\frac{D}{v}
\]

ดังนั้นช่วงเวลาระหว่างการชนของ $A$ กับ $B$ ครั้งแรกและครั้งที่สองคือ

\[
\Delta t=t_1+t_2
\]

\[
\boxed{\Delta t=\frac{2D}{v}}
\]
\end{examplebox}

% =========================================================
\subsection{ปี 2563 ข้อ 12: มวลไถลบนลิ่มที่เคลื่อนที่ได้}

\vspace{1.5em}

\IfFileExists{img/posn63q12.png}{
\begin{center}
    \includegraphics[width=1\linewidth]{img/posn63q12.png}
\end{center}
}{
\begin{tcolorbox}[title=ตำแหน่งภาพที่แนะนำ]
ให้ตัดภาพข้อสอบปี 2563 ข้อ 12 แล้วบันทึกเป็น

\[
\texttt{img/posn63q12.png}
\]
\end{tcolorbox}
}

\begin{examplebox}{เฉลยละเอียด}
มวล $m$ ไถลจากหยุดนิ่งลงจากความสูง $h$ บนลิ่มมวล $M$ พื้นทุกผิวลื่น เมื่อมวล $m$ หลุดจากปลายลิ่ม ความเร็วของทั้งสองวัตถุเป็นแนวราบ

ให้ความเร็วของ $m$ เทียบกับพื้นเป็น

\[
v
\]

ไปทางขวา และให้ลิ่มเคลื่อนที่ไปทางซ้ายด้วยอัตราเร็ว

\[
V
\]

ไม่มีแรงภายนอกในแนวราบ จึงอนุรักษ์โมเมนตัมแนวราบ

\[
mv-MV=0
\]

ดังนั้น

\[
V=\frac{m}{M}v
\]

พลังงานกลอนุรักษ์

\[
mgh=\frac{1}{2}mv^2+\frac{1}{2}MV^2
\]

แทน

\[
V=\frac{m}{M}v
\]

จะได้

\[
mgh=\frac{1}{2}mv^2+\frac{1}{2}M\left(\frac{m}{M}v\right)^2
\]

\[
mgh=\frac{1}{2}mv^2\left(1+\frac{m}{M}\right)
\]

จัดรูป

\[
v^2=\frac{2Mgh}{M+m}
\]

ดังนั้น

\[
\boxed{v=\sqrt{\frac{2Mgh}{M+m}}}
\]
\end{examplebox}

% =========================================================
\subsection{ปี 2563 ข้อ 13: การชนยืดหยุ่นแบบเยื้องศูนย์ของทรงกลมเท่ากัน}

\vspace{1.5em}

\IfFileExists{img/posn63q13.png}{
\begin{center}
    \includegraphics[width=1\linewidth]{img/posn63q13.png}
\end{center}
}{
\begin{tcolorbox}[title=ตำแหน่งภาพที่แนะนำ]
ให้ตัดภาพข้อสอบปี 2563 ข้อ 13 แล้วบันทึกเป็น

\[
\texttt{img/posn63q13.png}
\]
\end{tcolorbox}
}

\begin{examplebox}{เฉลยละเอียด}
ทรงกลม $A$ และ $B$ มีมวลเท่ากัน รัศมีเท่ากัน และผิวเกลี้ยงลื่น ขณะสัมผัสกัน แรงกระแทกเกิดตามแนวเส้นเชื่อมศูนย์กลางเท่านั้น

ให้แนวเส้นเชื่อมศูนย์กลางทำมุม $\theta$ กับทิศความเร็วเริ่มต้นของ $A$

จากเรขาคณิตของรูป ระยะเยื้องศูนย์คือ $h$ และระยะระหว่างศูนย์กลางขณะสัมผัสคือ $2R$

\[
\sin\theta=\frac{h}{2R}
\]

แยกความเร็วต้นของ $A$ เป็นสององค์ประกอบ

\[
v_{\parallel}=v\cos\theta
\]

ตามแนวเส้นเชื่อมศูนย์กลาง และ

\[
v_{\perp}=v\sin\theta
\]

ตั้งฉากกับแนวเส้นเชื่อมศูนย์กลาง

เพราะผิวเกลี้ยง แรงกระแทกไม่มีองค์ประกอบในแนวตั้งฉาก ดังนั้น $A$ ยังคงมีองค์ประกอบความเร็ว $v_{\perp}$ หลังชน ส่วนองค์ประกอบตามแนวเส้นเชื่อมศูนย์กลางถูกถ่ายโอนไปยัง $B$

จึงได้อัตราเร็วของ $A$ หลังชน

\[
v_A'=v\sin\theta
\]

แทน

\[
\sin\theta=\frac{h}{2R}
\]

ดังนั้น

\[
\boxed{v_A'=\frac{vh}{2R}}
\]
\end{examplebox}

% =========================================================
\subsection{ปี 2563 ข้อ 14: ทรงกระบอกสามท่อนและแรงตึงเชือก}

\vspace{1.5em}

\IfFileExists{img/posn63q14.png}{
\begin{center}
    \includegraphics[width=1\linewidth]{img/posn63q14.png}
\end{center}
}{
\begin{tcolorbox}[title=ตำแหน่งภาพที่แนะนำ]
ให้ตัดภาพข้อสอบปี 2563 ข้อ 14 แล้วบันทึกเป็น

\[
\texttt{img/posn63q14.png}
\]
\end{tcolorbox}
}

\begin{examplebox}{เฉลยละเอียด}
ทรงกระบอกทั้งสามมีรัศมีเท่ากัน จึงทำให้ศูนย์กลางของทรงกระบอกทั้งสามเป็นจุดยอดของสามเหลี่ยมด้านเท่า

ให้แรงปกติที่ทรงกระบอกล่างแต่ละท่อนกระทำต่อทรงกระบอกบนเป็น

\[
N
\]

แรงทั้งสองสมมาตรกันและทำมุม $60^\circ$ กับแนวราบ

สมดุลแรงแนวดิ่งของทรงกระบอกบนมวล $M$

\[
2N\sin60^\circ=Mg
\]

ดังนั้น

\[
N=\frac{Mg}{2\sin60^\circ}
\]

\[
N=\frac{Mg}{\sqrt{3}}
\]

สำหรับทรงกระบอกล่างแต่ละท่อน แรงจากทรงกระบอกบนมีองค์ประกอบแนวราบผลักออกด้านข้าง

\[
N\cos60^\circ
\]

เชือกต้องดึงกลับด้วยแรงตึงเท่ากัน

\[
T=N\cos60^\circ
\]

แทนค่า

\[
T=\frac{Mg}{\sqrt{3}}\left(\frac{1}{2}\right)
\]

ดังนั้น

\[
\boxed{T=\frac{Mg}{2\sqrt{3}}}
\]
\end{examplebox}

% =========================================================
\subsection{ปี 2563 ข้อ 20: พลังงานกลรวมของดาวเทียมในวงโคจรกลม}

\vspace{1.5em}

\IfFileExists{img/posn63q20.png}{
\begin{center}
    \includegraphics[width=1\linewidth]{img/posn63q20.png}
\end{center}
}{
\begin{tcolorbox}[title=ตำแหน่งภาพที่แนะนำ]
ให้ตัดภาพข้อสอบปี 2563 ข้อ 20 แล้วบันทึกเป็น

\[
\texttt{img/posn63q20.png}
\]
\end{tcolorbox}
}

\begin{examplebox}{เฉลยละเอียด}
ดาวเทียมมวล $m$ โคจรรอบโลกมวล $M$ เป็นวงกลมรัศมี $r$

แรงโน้มถ่วงทำหน้าที่เป็นแรงสู่ศูนย์กลาง

\[
\frac{GMm}{r^2}=\frac{mv^2}{r}
\]

ดังนั้น

\[
mv^2=\frac{GMm}{r}
\]

พลังงานจลน์คือ

\[
K=\frac{1}{2}mv^2
\]

\[
K=\frac{GMm}{2r}
\]

พลังงานศักย์โน้มถ่วงคือ

\[
U=-\frac{GMm}{r}
\]

ดังนั้น

\[
K=-\frac{1}{2}U
\]

พลังงานกลรวมคือ

\[
E=K+U
\]

\[
E=-\frac{1}{2}U+U
\]

\[
E=\frac{1}{2}U
\]

ดังนั้น

\[
\boxed{\frac{E}{U}=\frac{1}{2}}
\]
\end{examplebox}

% =========================================================
\subsection{ปี 2563 ข้อ 21: ลูกตุ้มกรวยและอัตราเร็วเชิงมุม}

\vspace{1.5em}

\IfFileExists{img/posn63q21.png}{
\begin{center}
    \includegraphics[width=1\linewidth]{img/posn63q21.png}
\end{center}
}{
\begin{tcolorbox}[title=ตำแหน่งภาพที่แนะนำ]
ให้ตัดภาพข้อสอบปี 2563 ข้อ 21 แล้วบันทึกเป็น

\[
\texttt{img/posn63q21.png}
\]
\end{tcolorbox}
}

\begin{examplebox}{เฉลยละเอียด}
ลูกตุ้มยาว $\ell$ หมุนเป็นวงกลมในระนาบระดับด้วยอัตราเร็วเชิงมุม $\omega$ และเชือกทำมุม $\theta$ กับแนวดิ่ง

รัศมีของวงกลมคือ

\[
r=\ell\sin\theta
\]

ให้แรงตึงเชือกเป็น $T$

สมดุลแนวดิ่ง

\[
T\cos\theta=mg
\]

แนวราบเป็นแรงสู่ศูนย์กลาง

\[
T\sin\theta=m\omega^2r
\]

แทน

\[
r=\ell\sin\theta
\]

จะได้

\[
T\sin\theta=m\omega^2\ell\sin\theta
\]

สำหรับกรณีที่ $\theta\neq 0$

\[
T=m\omega^2\ell
\]

แทนในสมการแนวดิ่ง

\[
m\omega^2\ell\cos\theta=mg
\]

ดังนั้น

\[
\cos\theta=\frac{g}{\omega^2\ell}
\]

จึงได้

\[
\boxed{\theta=\cos^{-1}\left(\frac{g}{\omega^2\ell}\right)}
\]
\end{examplebox}

% =========================================================
\subsection{ปี 2563 ข้อ 22: เพิ่มความเร็วของดาวเทียมให้หลุดจากวงโคจรกลม}

\vspace{1.5em}

\IfFileExists{img/posn63q22.png}{
\begin{center}
    \includegraphics[width=1\linewidth]{img/posn63q22.png}
\end{center}
}{
\begin{tcolorbox}[title=ตำแหน่งภาพที่แนะนำ]
ให้ตัดภาพข้อสอบปี 2563 ข้อ 22 แล้วบันทึกเป็น

\[
\texttt{img/posn63q22.png}
\]
\end{tcolorbox}
}

\begin{examplebox}{เฉลยละเอียด}
ดาวเทียมอยู่ในวงโคจรกลมรัศมี $r$ ด้วยอัตราเร็ว $v_0$

จากแรงสู่ศูนย์กลาง

\[
\frac{GMm}{r^2}=\frac{mv_0^2}{r}
\]

ดังนั้น

\[
v_0^2=\frac{GM}{r}
\]

อัตราเร็วหลุดพ้น $v_e$ หาได้จากเงื่อนไขพลังงานกลรวมเท่ากับศูนย์

\[
\frac{1}{2}mv_e^2-\frac{GMm}{r}=0
\]

ดังนั้น

\[
v_e^2=\frac{2GM}{r}
\]

เปรียบเทียบกับ

\[
v_0^2=\frac{GM}{r}
\]

จะได้

\[
v_e^2=2v_0^2
\]

ดังนั้น

\[
\boxed{v_e=\sqrt{2}\,v_0}
\]
\end{examplebox}

% =========================================================
\subsection{ปี 2563 ข้อ 23: ดาวเทียมในวงโคจรวงรีและโมเมนตัมเชิงมุม}

\vspace{1.5em}

\IfFileExists{img/posn63q23.png}{
\begin{center}
    \includegraphics[width=1\linewidth]{img/posn63q23.png}
\end{center}
}{
\begin{tcolorbox}[title=ตำแหน่งภาพที่แนะนำ]
ให้ตัดภาพข้อสอบปี 2563 ข้อ 23 แล้วบันทึกเป็น

\[
\texttt{img/posn63q23.png}
\]
\end{tcolorbox}
}

\begin{examplebox}{เฉลยละเอียด}
ดาวเทียมมวล $m$ โคจรรอบโลกเป็นวงรี ที่จุดใกล้สุดและจุดไกลสุด ความเร็วตั้งฉากกับแนวรัศมี

แรงโน้มถ่วงเป็นแรงเข้าสู่ศูนย์กลาง จึงให้ทอร์กรอบศูนย์กลางโลกเป็นศูนย์

\[
\tau=0
\]

ดังนั้นโมเมนตัมเชิงมุมอนุรักษ์

\[
L_1=L_2
\]

ที่จุดใกล้สุด

\[
L_1=mr_1v_1
\]

ที่จุดไกลสุด

\[
L_2=mr_2v_2
\]

ดังนั้น

\[
mr_1v_1=mr_2v_2
\]

ตัด $m$

\[
r_1v_1=r_2v_2
\]

จึงได้

\[
\boxed{v_2=\frac{r_1}{r_2}v_1}
\]
\end{examplebox}

% =========================================================
\subsection{ปี 2563 ข้อ 26: ระดับของเหลวในท่อเมื่อความดันต่างกัน}

\vspace{1.5em}

\IfFileExists{img/posn63q26.png}{
\begin{center}
    \includegraphics[width=1\linewidth]{img/posn63q26.png}
\end{center}
}{
\begin{tcolorbox}[title=ตำแหน่งภาพที่แนะนำ]
ให้ตัดภาพข้อสอบปี 2563 ข้อ 26 แล้วบันทึกเป็น

\[
\texttt{img/posn63q26.png}
\]
\end{tcolorbox}
}

\begin{examplebox}{เฉลยละเอียด}
อากาศในปลายท่อมีความดัน

\[
P_i
\]

ส่วนผิวของเหลวนอกท่อสัมผัสบรรยากาศที่ความดัน

\[
P_0
\]

ผิวของเหลวในท่อสูงกว่าผิวนอกเป็นระยะ

\[
h
\]

พิจารณาจุดสองจุดที่อยู่ในของเหลวระดับเดียวกัน จุดหนึ่งอยู่นอกท่อที่ระดับผิวนอก และอีกจุดอยู่ในท่อ

ความดันนอกท่อคือ

\[
P_0
\]

ความดันในท่อที่ระดับเดียวกันคือ

\[
P_i+\rho gh
\]

ความดันที่ระดับเดียวกันในของเหลวชนิดเดียวกันต้องเท่ากัน

\[
P_i+\rho gh=P_0
\]

ดังนั้น

\[
\boxed{h=\frac{P_0-P_i}{\rho g}}
\]
\end{examplebox}

% =========================================================
\subsection{ปี 2563 ข้อ 27: ทิศของความเร็วโพรเจคไทล์ ณ เวลาใด ๆ}

\vspace{1.5em}

\IfFileExists{img/posn63q27.png}{
\begin{center}
    \includegraphics[width=1\linewidth]{img/posn63q27.png}
\end{center}
}{
\begin{tcolorbox}[title=ตำแหน่งภาพที่แนะนำ]
ให้ตัดภาพข้อสอบปี 2563 ข้อ 27 แล้วบันทึกเป็น

\[
\texttt{img/posn63q27.png}
\]
\end{tcolorbox}
}

\begin{examplebox}{เฉลยละเอียด}
ยิงโพรเจคไทล์ด้วยความเร็วต้น $u$ ทำมุม $\theta_0$ กับแกน $+x$

ความเร็วแนวราบคงตัว

\[
v_x=u\cos\theta_0
\]

ความเร็วแนวดิ่งที่เวลา $t$ คือ

\[
v_y=u\sin\theta_0-gt
\]

ให้เวกเตอร์ความเร็วทำมุม $\phi$ กับแกน $+x$

\[
\tan\phi=\frac{v_y}{v_x}
\]

แทนค่า

\[
\tan\phi=
\frac{u\sin\theta_0-gt}{u\cos\theta_0}
\]

ดังนั้น

\[
\boxed{
\tan\phi=
\frac{u\sin\theta_0-gt}{u\cos\theta_0}
}
\]
\end{examplebox}

% =========================================================
\subsection{ปี 2563 ข้อ 28: ยิงโพรเจคไทล์สองลูกต่างตำแหน่งและต่างเวลา}

\vspace{1.5em}

\IfFileExists{img/posn63q28.png}{
\begin{center}
    \includegraphics[width=1\linewidth]{img/posn63q28.png}
\end{center}
}{
\begin{tcolorbox}[title=ตำแหน่งภาพที่แนะนำ]
ให้ตัดภาพข้อสอบปี 2563 ข้อ 28 แล้วบันทึกเป็น

\[
\texttt{img/posn63q28.png}
\]
\end{tcolorbox}
}

\begin{examplebox}{เฉลยละเอียด}
จุด $A$ และ $B$ อยู่ห่างกันเป็นระยะ

\[
D
\]

ยิงลูกแรกจาก $A$ ด้วยความเร็วต้น $u$ และมุม $\theta_0$ จากนั้นยิงลูกที่สองจาก $B$ ด้วยความเร็วต้นและมุมเท่ากัน ในจังหวะที่ลูกแรกอยู่เหนือจุด $B$ พอดี

ความเร็วแนวราบของทั้งสองลูกคือ

\[
u_x=u\cos\theta_0
\]

ลูกแรกใช้เวลาเดินทางจาก $A$ ไปถึงแนวดิ่งเหนือ $B$

\[
t_0=\frac{D}{u\cos\theta_0}
\]

ให้ $t$ เป็นเวลาหลังจากยิงลูกที่สอง

เพราะทั้งสองลูกมีความเร็วแนวราบเท่ากัน หลังจากยิงลูกที่สอง ตำแหน่งแนวราบของทั้งสองลูกจะตรงกันตลอดเวลา

พิจารณาแกน $y$

ตำแหน่งแนวดิ่งของลูกแรกคือ

\[
y_1=u\sin\theta_0(t+t_0)-\frac{1}{2}g(t+t_0)^2
\]

ตำแหน่งแนวดิ่งของลูกที่สองคือ

\[
y_2=u\sin\theta_0t-\frac{1}{2}gt^2
\]

เมื่อชนกัน

\[
y_1=y_2
\]

ลบสมการและตัด $t_0\neq 0$

\[
u\sin\theta_0-g t-\frac{1}{2}g t_0=0
\]

ดังนั้น

\[
t=\frac{u\sin\theta_0}{g}-\frac{t_0}{2}
\]

ตำแหน่งที่ชนวัดจากจุด $A$ คือ

\[
x=D+u\cos\theta_0\,t
\]

แทนค่า $t$

\[
x=D+u\cos\theta_0
\left(
\frac{u\sin\theta_0}{g}-\frac{t_0}{2}
\right)
\]

และใช้

\[
t_0=\frac{D}{u\cos\theta_0}
\]

จะได้

\[
x=D+\frac{u^2\sin\theta_0\cos\theta_0}{g}-\frac{D}{2}
\]

ดังนั้น

\[
\boxed{
x=\frac{D}{2}+\frac{u^2\sin\theta_0\cos\theta_0}{g}
}
\]

หรือใช้ $\sin2\theta_0=2\sin\theta_0\cos\theta_0$

\[
\boxed{
x=\frac{D}{2}+\frac{u^2\sin2\theta_0}{2g}
}
\]
\end{examplebox}

% =========================================================
\subsection{ปี 2563 ข้อ 29: บอลลูนอากาศร้อนและแรงยกลัพธ์ \textnormal{(ข้อเสริมเชื่อมกับแก๊สอุดมคติ)}}

\vspace{1.5em}

\IfFileExists{img/posn63q29.png}{
\begin{center}
    \includegraphics[width=1\linewidth]{img/posn63q29.png}
\end{center}
}{
\begin{tcolorbox}[title=ตำแหน่งภาพที่แนะนำ]
ให้ตัดภาพข้อสอบปี 2563 ข้อ 29 แล้วบันทึกเป็น

\[
\texttt{img/posn63q29.png}
\]
\end{tcolorbox}
}

\begin{examplebox}{เฉลยละเอียด}
โจทย์กำหนดให้มวลโมลาร์ของอากาศเป็น $M$ ความดันภายในและภายนอกบอลลูนเท่ากันเป็น $P$ ปริมาตรบอลลูนเป็น $V$ อุณหภูมิอากาศภายในเป็น $T_1$ และอุณหภูมิอากาศภายนอกเป็น $T_0$

จากกฎแก๊สอุดมคติ

\[
PV=nRT
\]

และเพราะมวลแก๊สเท่ากับ

\[
m=nM
\]

ความหนาแน่นของแก๊สคือ

\[
\rho=\frac{m}{V}
\]

จึงได้

\[
\rho=\frac{PM}{RT}
\]

ความหนาแน่นอากาศภายนอกคือ

\[
\rho_0=\frac{PM}{RT_0}
\]

ความหนาแน่นอากาศร้อนภายในคือ

\[
\rho_1=\frac{PM}{RT_1}
\]

แรงพยุงมีขนาด

\[
B=\rho_0Vg
\]

น้ำหนักของอากาศร้อนภายในบอลลูนคือ

\[
W=\rho_1Vg
\]

ดังนั้นแรงยกลัพธ์คือ

\[
F_{\text{lift}}=B-W
\]

\[
F_{\text{lift}}=(\rho_0-\rho_1)Vg
\]

แทนค่าความหนาแน่น

\[
\boxed{
F_{\text{lift}}
=\frac{PMVg}{R}
\left(
\frac{1}{T_0}-\frac{1}{T_1}
\right)
}
\]
\end{examplebox}

% =========================================================
\subsection{ปี 2563 ข้อ 31: มวลไถลหลุดจากผิวทรงกลม}

\vspace{1.5em}

\IfFileExists{img/posn63q31.png}{
\begin{center}
    \includegraphics[width=1\linewidth]{img/posn63q31.png}
\end{center}
}{
\begin{tcolorbox}[title=ตำแหน่งภาพที่แนะนำ]
ให้ตัดภาพข้อสอบปี 2563 ข้อ 31 แล้วบันทึกเป็น

\[
\texttt{img/posn63q31.png}
\]
\end{tcolorbox}
}

\begin{examplebox}{เฉลยละเอียด}
มวล $m$ เริ่มจากหยุดนิ่งที่ยอดผิวทรงกลมรัศมี $R$ และไถลลงบนผิวลื่น

ให้มวลอยู่ที่ตำแหน่งซึ่งแนวรัศมีทำมุม $\theta$ กับแนวดิ่ง

มวลลดระดับลงเป็นระยะ

\[
\Delta h=R-R\cos\theta
\]

\[
\Delta h=R(1-\cos\theta)
\]

ใช้การอนุรักษ์พลังงาน

\[
mg\Delta h=\frac{1}{2}mv^2
\]

ดังนั้น

\[
mgR(1-\cos\theta)=\frac{1}{2}mv^2
\]

\[
v^2=2gR(1-\cos\theta)
\]

พิจารณาแรงในแนวเข้าสู่ศูนย์กลาง

\[
mg\cos\theta-N=\frac{mv^2}{R}
\]

เมื่อมวลกำลังหลุดจากผิวทรงกลม แรงปกติเป็นศูนย์

\[
N=0
\]

จึงได้

\[
mg\cos\theta=\frac{mv^2}{R}
\]

\[
v^2=gR\cos\theta
\]

นำไปเทียบกับสมการพลังงาน

\[
gR\cos\theta=2gR(1-\cos\theta)
\]

\[
\cos\theta=2-2\cos\theta
\]

\[
3\cos\theta=2
\]

\[
\cos\theta=\frac{2}{3}
\]

ดังนั้น

\[
v^2=gR\left(\frac{2}{3}\right)
\]

จึงได้

\[
\boxed{v=\sqrt{\frac{2gR}{3}}}
\]
\end{examplebox}

% =========================================================
\section{โจทย์ประยุกต์รวมกลศาสตร์ที่แต่งเพิ่ม}

\begin{exercisebox}{ชุดโจทย์ประยุกต์รวมกลศาสตร์}
\begin{enumerate}
    \item กล่องมวล $m$ เริ่มจากหยุดนิ่งบนพื้นเอียงลื่นสูง $h$ แล้วไถลลงมาชนสปริงค่าคงตัว $k$ ที่ปลายล่าง จงหาระยะอัดมากที่สุดของสปริง ถ้าระดับของสปริงอยู่ที่ฐานพื้นเอียง

    \item วัตถุมวล $m$ ผูกเชือกยาว $R$ แกว่งในแนวดิ่ง ถ้าที่ตำแหน่งล่างสุดมีอัตราเร็ว $v$ จงหาแรงตึงเชือกที่ตำแหน่งล่างสุด และหาเงื่อนไขขั้นต่ำของอัตราเร็วที่ตำแหน่งบนสุดเพื่อให้เชือกไม่หย่อน

    \item ลูกบอลมวล $m$ พุ่งชนกำแพงด้วยความเร็ว $u$ ในแนวตั้งฉากกับกำแพง แล้วสะท้อนกลับด้วยความเร็ว $eu$ โดย $0<e<1$ ถ้าเวลาสัมผัสกำแพงคือ $\Delta t$ จงหาแรงเฉลี่ยจากกำแพง

    \item วัตถุทรงกระบอกตันมวล $M$ รัศมี $R$ กลิ้งโดยไม่ไถลลงพื้นเอียงมุม $\theta$ จงหาความเร่งของศูนย์กลางมวล โดยใช้ $I=\frac{1}{2}MR^2$

    \item วัตถุความหนาแน่น $\rho_o$ ลอยในของเหลวความหนาแน่น $\rho_f$ โดยมีปริมาตรรวม $V$ ถ้ากดวัตถุให้จมทั้งก้อนแล้วปล่อย จงหาแรงลัพธ์เริ่มต้นที่กระทำต่อวัตถุ
\end{enumerate}
\end{exercisebox}

\section{เฉลยโจทย์ประยุกต์รวมกลศาสตร์}

\begin{enumerate}
    \item
    ใช้การอนุรักษ์พลังงานกล ถ้าฐานพื้นเอียงเป็นระดับอ้างอิง พลังงานศักย์เริ่มต้นคือ
    \[
    mgh
    \]
    ที่อัดสปริงมากที่สุด ความเร็วเป็นศูนย์ และพลังงานทั้งหมดอยู่ในสปริง
    \[
    \frac{1}{2}kx^2=mgh
    \]
    ดังนั้น
    \[
    \boxed{x=\sqrt{\frac{2mgh}{k}}}
    \]

    \item
    ที่ตำแหน่งล่างสุด แรงสู่ศูนย์กลางชี้ขึ้น
    \[
    T-mg=\frac{mv^2}{R}
    \]
    ดังนั้น
    \[
    \boxed{T=mg+\frac{mv^2}{R}}
    \]
    ที่ตำแหน่งบนสุด เงื่อนไขเชือกไม่หย่อนคือแรงตึงน้อยสุดเท่ากับศูนย์
    \[
    mg=\frac{mv_{\min}^2}{R}
    \]
    ดังนั้น
    \[
    \boxed{v_{\min}=\sqrt{gR}}
    \]

    \item
    เลือกทิศเข้าหากำแพงเป็นบวก ความเร็วต้นคือ
    \[
    u
    \]
    ความเร็วหลังสะท้อนคือ
    \[
    -eu
    \]
    การเปลี่ยนโมเมนตัมของลูกบอลคือ
    \[
    \Delta p=m(-eu)-mu
    \]
    \[
    \Delta p=-m(1+e)u
    \]
    ขนาดแรงเฉลี่ยคือ
    \[
    \boxed{F_{\text{avg}}=\frac{m(1+e)u}{\Delta t}}
    \]

    \item
    สำหรับการกลิ้งโดยไม่ไถล
    \[
    a=\frac{g\sin\theta}{1+\frac{I}{MR^2}}
    \]
    แทน
    \[
    I=\frac{1}{2}MR^2
    \]
    จะได้
    \[
    a=\frac{g\sin\theta}{1+\frac{1}{2}}
    \]
    \[
    \boxed{a=\frac{2}{3}g\sin\theta}
    \]

    \item
    เมื่อวัตถุจมทั้งก้อน แรงพยุงคือ
    \[
    B=\rho_fgV
    \]
    น้ำหนักคือ
    \[
    W=\rho_ogV
    \]
    ดังนั้นแรงลัพธ์ขึ้นคือ
    \[
    F_{\text{net}}=B-W
    \]
    \[
    \boxed{F_{\text{net}}=(\rho_f-\rho_o)gV}
    \]
\end{enumerate}
% =========================================================
\section{สรุปเครื่องมือสำหรับโจทย์รวมกลศาสตร์}

\begin{tcolorbox}[title=สูตรหลักที่ใช้บ่อยในข้อสอบรวม]
การเคลื่อนที่ด้วยความเร่งคงตัว

\[
v_f=v_i+at
\]

\[
\Delta x=v_it+\frac{1}{2}at^2
\]

\[
v_f^2=v_i^2+2a\Delta x
\]

กฎของนิวตัน

\[
\sum \vec{F}=m\vec{a}
\]

งานและพลังงาน

\[
W_{\text{net}}=\Delta K
\]

\[
K_i+U_i+W_{\text{nonconservative}}=K_f+U_f
\]

โมเมนตัมและอิมพัลส์

\[
\vec{J}=\Delta \vec{p}
\]

\[
\sum \vec{p}_i=\sum \vec{p}_f
\]

ทอร์กและการหมุน

\[
\sum \tau=I\alpha
\]

\[
\sum F_x=0,\quad \sum F_y=0,\quad \sum \tau=0
\]

ของไหล

\[
P=P_0+\rho gh
\]

\[
B=\rho_fgV_{\text{disp}}
\]

\[
A_1v_1=A_2v_2
\]

\[
P+\frac{1}{2}\rho v^2+\rho gy=\text{constant}
\]
\end{tcolorbox}